%
%     Mathematical relationship(s) between CD matrix and CRPIXi.
%
% Rob Siverd
% Created:       2026-03-17
% Last modified: 2026-03-24
%---------------------------------------------------------------------
%*********************************************************************
%---------------------------------------------------------------------

\documentclass{article}

\usepackage{color}
\usepackage{ulem}
\usepackage{cancel}
\usepackage{amsmath}

\pdfpagewidth 8.5in
\pdfpageheight 11in
\setlength\topmargin{0.2in}
\setlength\headheight{0in}
\setlength\textheight{8in}
\setlength\textwidth{6in}
\setlength\oddsidemargin{0in}
\setlength\evensidemargin{0in}

\begin{document}

% Units:
\newcommand{\Msun}{M_{sun}}
\newcommand{\Lsun}{L_{sun}}
\newcommand{\kg}{\mathrm{kg}}
\newcommand{\K}{\mathrm{K}}
\newcommand{\nm}{\mathrm{nm}}
\newcommand{\cm}{\mathrm{cm}}
\newcommand{\km}{\mathrm{km}}
\newcommand{\erg}{\mathrm{erg}}
\newcommand{\cms}{\mathrm{cm/s}}
\newcommand{\kms}{\mathrm{km/s}}
\newcommand{\mph}{\mathrm{mph}}
\newcommand{\mpc}{\mathrm{Mpc}}
\newcommand{\mi}{\mathrm{mi}}
\newcommand{\au}{\mathrm{AU}}
\newcommand{\Sec}{\mathrm{s}}
\newcommand{\Min}{\mathrm{min}}
\newcommand{\Hrs}{\mathrm{hours}}
\newcommand{\Hr}{\mathrm{hr}}
\newcommand{\ly}{\mathrm{ly}}
\newcommand{\yr}{\mathrm{yr}}
\newcommand{\pc}{\mathrm{pc}}
\newcommand{\strs}{\mathrm{stars}}
\newcommand{\partdiv}{------------------------------------------}
\newcommand{\divider}{\partdiv \partdiv \partdiv}

%\newcommand{\crpixone}{\mathrm{CRPIX1}}
%\newcommand{\crpixi}[1]{\mathrm{CRPIX#1}}
\newcommand{\cpx}[1]{\mathrm{CRPIX#1}}
\newcommand{\pcm}[1]{\mathrm{PC}_{#1}}
\newcommand{\cdm}[1]{\mathrm{CD}_{#1}}

\newcommand{\xccd}{\mathrm{X}_{\mathrm{DET}}}
\newcommand{\yccd}{\mathrm{Y}_{\mathrm{DET}}}
\newcommand{\xfoc}{\mathrm{X}_{\mathrm{FOC}}}
\newcommand{\yfoc}{\mathrm{Y}_{\mathrm{FOC}}}

\newcommand{\COS}[1]{\mathrm{cos}(#1)}
\newcommand{\SIN}[1]{\mathrm{sin}(#1)}
\newcommand{\RMCOS}{\COS{\theta}}
\newcommand{\RMSIN}{\SIN{\theta}}

% Cancel text in red:
\renewcommand{\CancelColor}{\thicklines\color{red}}


I am using this document to stumble my way through the transformation
algebra needed to related CRPIXi values among different WIRCam sensors.\\

In what follows, we use these symbols:
\begin{flalign*}
   \xccd &= \text{detector X pixel coordinate} & \\
   \yccd &= \text{detector Y pixel coordinate} & \\
   \xfoc &= \text{X coordinate in focal plane} & \\
   \yfoc &= \text{Y coordinate in focal plane} &
\end{flalign*}

The FITS WCS papers give formulae:
\begin{flalign*}
   x_i  =  s_i * &\sum \left[       m_{ij} * (p_j - r_j)  \right] & \\
        =        &\sum \left[ s_i * m_{ij} * (p_j - r_j)  \right] &
\end{flalign*}

where:
\begin{flalign*}
   s_i          &== \textrm{CDELTia}   & \\
         m_{ij} &== \textrm{PCi\_ja}   & \\
   s_i * m_{ij} &== \textrm{CDi\_ja}   & 
\end{flalign*}\\

The last line above can be written as follows in matrix notation:
\[
   \begin{pmatrix}
      \cdm{11} & \cdm{12} \\
      \cdm{21} & \cdm{22} \\
   \end{pmatrix}
   =
   \begin{pmatrix}
      s_1 &   0 \\
        0 & s_2 \\
   \end{pmatrix}
   \cdot
   \begin{pmatrix}
      \pcm{11} & \pcm{12} \\
      \pcm{21} & \pcm{22} \\
   \end{pmatrix}
\]\\

\hrule
\vspace{2mm}

A reminder about VECTOR (active) and FRAME (passive) rotation matrix conventions:
\begin{align}
   R_{zV}(\theta) &=
   \begin{pmatrix}
      \RMCOS & -\RMSIN \\
      \RMSIN & \;\;\RMCOS
   \end{pmatrix} \\
   R_{zF}(\theta) &=
   \begin{pmatrix}
      \;\;\RMCOS & \RMSIN \\
      -\RMSIN & \RMCOS
   \end{pmatrix}
\end{align}



\pagebreak

Definitions from the WCS papers spelled out (PC formalism):
%\[
\begin{flalign}
%  \xi  &= s_1 * \textrm{PC}_{11} * (p_x - \cpx{1}) + s1 * PC1_2 * (p_y - \cpx{2}) \\
%  \eta &= s_2 * \textrm{PC}_{21} * (p_x - \cpx{1}) + s2 * PC2_2 * (p_y - \cpx{2})
   \xfoc &= s_1 \cdot \pcm{11} \cdot (\xccd - \cpx{1}) \;\; + s1 \cdot \pcm{12} \cdot (\yccd - \cpx{2}) \\
   \yfoc &= s_2 \cdot \pcm{21} \cdot (\xccd - \cpx{1}) \;\; + s2 \cdot \pcm{22} \cdot (\yccd - \cpx{2}) 
\end{flalign}
\vspace{-7mm}
\begin{flalign}
   \label{eq:xfoc} \xfoc &= \cdm{11} \cdot (\xccd - \cpx{1}) \;\; + \cdm{12} \cdot (\yccd - \cpx{2}) \\
   \label{eq:yfoc} \yfoc &= \cdm{21} \cdot (\xccd - \cpx{1}) \;\; + \cdm{22} \cdot (\yccd - \cpx{2})
\end{flalign}\\

We can put these equations into matrix form using homogeneous 
coordinates and augmented matrices. For translation:
\[
   \begin{pmatrix}
      1 & 0 & -\cpx{1} \\
      0 & 1 & -\cpx{2} \\
      0 & 0 & 1
   \end{pmatrix}
   \cdot
   \begin{pmatrix}
      \xccd \\
      \yccd \\
      1
   \end{pmatrix}
   =
   \begin{pmatrix}
      \xccd \;-\; \cpx{1} \\
      \yccd \;-\; \cpx{2} \\
      1
   \end{pmatrix}
   \rightarrow
   \begin{pmatrix}
      A \\
      B \\
      1
   \end{pmatrix}
\]

For the CD matrix components (AFTER translation), we do:
\[
   \begin{pmatrix}
      \cdm{11} & \cdm{12} & 0 \\
      \cdm{21} & \cdm{22} & 0 \\
         0     &    0     & 1
   \end{pmatrix}
   \cdot
   \begin{pmatrix}
      A \\
      B \\
      1
   \end{pmatrix}
   =
   \begin{pmatrix}
      \cdm{11} \cdot A \; + \; \cdm{12} \cdot B \\
      \cdm{21} \cdot A \; + \; \cdm{22} \cdot B \\
      1
   \end{pmatrix}
\]

Combining these, we get:
\begin{align}
   \begin{pmatrix}
      \xfoc \\
      \yfoc \\
      1
   \end{pmatrix}
   &=
   \begin{pmatrix}
      \cdm{11} & \cdm{12} & 0 \\
      \cdm{21} & \cdm{22} & 0 \\
         0     &    0     & 1
   \end{pmatrix}
   \cdot
   \begin{pmatrix}
      1 & 0 & -\cpx{1} \\
      0 & 1 & -\cpx{2} \\
      0 & 0 & 1
   \end{pmatrix}
   \cdot
   \begin{pmatrix}
      \xccd \\
      \yccd \\
      1
   \end{pmatrix} \\
   &=
   \begin{pmatrix}
      \cdm{11} & \cdm{12} & -(\cdm{11} \cdot \cpx{1} + \cdm{12} \cdot \cpx{2}) \\
      \cdm{21} & \cdm{22} & -(\cdm{21} \cdot \cpx{1} + \cdm{22} \cdot \cpx{2}) \\
         0     &    0     & 1
   \end{pmatrix}
   \cdot
   \begin{pmatrix}
      \xccd \\
      \yccd \\
      1
   \end{pmatrix}
\end{align}

% -----------------------------------------------------------------------

\hrule
\vspace{2mm}
\hrule
\vspace{5mm}

\pagebreak

Next we determine how the various CDij and CRPIXn are related between the sensors.
We showed in Eqns \ref{eq:xfoc}, \ref{eq:yfoc} how the x,y focal plane coordinates
depend on sensor x,y, CD matrix values, and CRPIXn. Recall from above that
(where subscript q now implies sensor):
these equations illustrate the connections between mosaic position and individual
sensor parameters (where subscript q implies sensor). Recall:
\begin{align*}
   \xfoc &= \cdm{11,q} \cdot (\xccd - \cpx{1}_q) \;\; + \cdm{12,q} \cdot (\yccd - \cpx{2}_q) \\
   \yfoc &= \cdm{21,q} \cdot (\xccd - \cpx{1}_q) \;\; + \cdm{22,q} \cdot (\yccd - \cpx{2}_q)
\end{align*}
%Let $A == \cpx{1}$ and $B = \cpx{2}$ for brevity. Recall:
%\begin{align*}
%   \xfoc &= \cdm{11,q} \cdot (\xccd - A_q) \;\; + \cdm{12,q} \cdot (\yccd - B_q) \\
%   \yfoc &= \cdm{21,q} \cdot (\xccd - A_q) \;\; + \cdm{22,q} \cdot (\yccd - B_q)
%\end{align*}
%.
The total differentials below (taken holding detector coordinates and CD 
matrix entries constant) describe the connection between mosaic position in 
the focal plane and sensor parameters:
\begin{align}
   d(\xfoc) &= \frac{\partial (\xfoc)}{\partial (\cpx{1}_q)} d(\cpx{1}_q) 
               + \frac{\partial (\xfoc)}{\partial (\cpx{2}_q)} d(\cpx{2}_q) \\
   d(\yfoc) &= \frac{\partial (\yfoc)}{\partial (\cpx{1}_q)} d(\cpx{1}_q) 
               + \frac{\partial (\yfoc)}{\partial (\cpx{2}_q)} d(\cpx{2}_q)
\end{align}
%\begin{align}
%   d(\xfoc) &= \frac{\partial (\xfoc)}{\partial (A_q)} d(A_q) 
%               + \frac{\partial (\xfoc)}{\partial (B_q)} d(B_q) \\
%   d(\yfoc) &= \frac{\partial (\yfoc)}{\partial (A_q)} d(A_q) 
%               + \frac{\partial (\yfoc)}{\partial (B_q)} d(B_q)
%\end{align}
where the partial differentials are:
\begin{align*}
   \frac{\partial (\xfoc)}{\partial (\cpx{1}_q)} &= -\cdm{11,q} &
   \frac{\partial (\xfoc)}{\partial (\cpx{2}_q)} &= -\cdm{12,q} \\
   \frac{\partial (\yfoc)}{\partial (\cpx{1}_q)} &= -\cdm{21,q} &
   \frac{\partial (\yfoc)}{\partial (\cpx{2}_q)} &= -\cdm{22,q} \\
\end{align*}
%\begin{align}
%   \frac{\partial (\xfoc)}{\partial (A_q)} &= 
%              -\left[\cdm{11,q} \cdot \cpx{1} + \cdm{12,q} \cdot \cpx{2} \right] \\
%   \frac{\partial (\yfoc)}{\partial (\cpx{1}_q)} &= 
%              -\left[\cdm{21,q} \cdot \cpx{1} + \cdm{22,q} \cdot \cpx{2} \right]
%\end{align}

The total differentials thus become:
\begin{align}
   \Delta \xfoc &= -\cdm{11,q} \cdot \Delta \cpx{1}_q - \cdm{12,q} \cdot \Delta \cpx{2}_q \\
   \Delta \yfoc &= -\cdm{21,q} \cdot \Delta \cpx{1}_q - \cdm{22,q} \cdot \Delta \cpx{2}_q
\end{align}

\hrule
\vspace{2mm}
\hrule
\vspace{5mm}

\pagebreak

New approach: we put all of the sensors directly onto the NE pixel grid.
\\

In matrix form, CD (1st) and translation (2nd) combine as follows:
\begin{align*}
   \begin{pmatrix}
      1 & 0 & T_x \\
      0 & 1 & T_y \\
      0 & 0 & 1
   \end{pmatrix}
   \begin{pmatrix}
      \cdm{11} & \cdm{12} & 0 \\
      \cdm{21} & \cdm{22} & 0 \\
         0     &    0     & 1
   \end{pmatrix}
   =
   \begin{pmatrix}
      \cdm{11} & \cdm{12} & T_x \\
      \cdm{21} & \cdm{22} & T_y \\
         0     &    0     & 1
   \end{pmatrix}
\end{align*}

Rotation BEFORE scale (not preferred) goes as follows:
\begin{align*}
   \begin{pmatrix}
         s_1 &  0   \\
         0   & s_2 
   \end{pmatrix}
   \begin{pmatrix}
      \RMCOS   &    -\RMSIN   \\
      \RMSIN   & \;\;\RMCOS 
   \end{pmatrix}
   =
   \begin{pmatrix}
      s_1 \, \RMCOS   &     -s_1 \, \RMSIN   \\
      s_2 \, \RMSIN   & \;\; s_2 \, \RMCOS 
   \end{pmatrix}
\end{align*}

Rotation AFTER scale (preferred for SRT) goes as follows:
\begin{align*}
   \begin{pmatrix}
      \RMCOS   &    -\RMSIN   \\
      \RMSIN   & \;\;\RMCOS 
   \end{pmatrix}
   \begin{pmatrix}
         s_1 &  0   \\
         0   & s_2 
   \end{pmatrix}
   =
   \begin{pmatrix}
      s_1 \, \RMCOS   &     -s_2 \, \RMSIN   \\
      s_1 \, \RMSIN   & \;\; s_2 \, \RMCOS 
   \end{pmatrix}
\end{align*}

The full SRT transformation series in augmented (affine
transformation) matrix form:
\begin{align*}
   \begin{pmatrix}
      1 & 0 & T_x \\
      0 & 1 & T_y \\
      0 & 0 & 1
   \end{pmatrix}
   \begin{pmatrix}
      \RMCOS   &    -\RMSIN & 0  \\
      \RMSIN   & \;\;\RMCOS & 0  \\
           0   &          0 & 1
   \end{pmatrix}
   \begin{pmatrix}
      s_1 &   0  & 0  \\
        0 & s_2  & 0  \\
        0 &   0  & 1 
   \end{pmatrix}
   =
   \begin{pmatrix}
      s_1 \, \RMCOS   &     -s_2 \, \RMSIN  & T_x \\
      s_1 \, \RMSIN   & \;\; s_2 \, \RMCOS  & T_y \\
                  0   &                  0  &   1
   \end{pmatrix}
\end{align*}

The SRT order is nice because (a) any scale changes can be interpreted as
tip/tilt in the detector plane along its native pixel axes, (b) detector
rotation, shear, and/or skew all occur in the NE sensor pixel plane, and 
(c) the SRT transformation is fairly standard in computer graphics and
numerous tools and resources exist for its interpretation as a result.\\

\hrule
\vspace{2mm}
\hrule
\vspace{5mm}

%A reminder about VECTOR (active) and FRAME (passive) rotation matrix conventions:
%\begin{align}
%   R_{zV}(\theta) &=
%   \begin{pmatrix}
%      \RMCOS & -\RMSIN \\
%      \RMSIN & \;\;\RMCOS
%   \end{pmatrix} \\
%   R_{zF}(\theta) &=
%   \begin{pmatrix}
%      \;\;\RMCOS & \RMSIN \\
%      -\RMSIN & \RMCOS
%   \end{pmatrix}
%\end{align}

\pagebreak

\section{Mosaic Transformation Models}

The following sections describe different possible transformations from the
other (NW, SE, SW) sensors to the NE sensor. They are presented in increasing
order of complication (parameter count). A wise course of action is fitting
more than one model to compare results. Since these models are nested (or
very close to it) we should be able to formally determine whether the more
complicated models are justified statistically.

\subsection{Rotation and Translation}

The simplest case is assuming that all detectors are coplanar with matched scale
(they differ by only a rotation and a translation). In the following, $\theta$ is
the counter-clockwise rotation angle of detector Q axes w.r.t. NE detector axes. 
This approach uses \textbf{3 parameters} (rotation $\theta$ and translations 
$O_x$ and $O_y$) for each of the (NW, SE, and SW) sensors for 
\textbf{a total of 9 parameters}.
\begin{align*}
%\[
   \begin{pmatrix}
      x \\
      y
   \end{pmatrix}_{NE}
   &=
   R_{zV}(\theta)
   \begin{pmatrix}
      x \\
      y
   \end{pmatrix}_{Q}
   +
   \begin{pmatrix}
      O_x \\
      O_y
   \end{pmatrix}_{Q}
   =
   R_{zF}(-\theta)
   \begin{pmatrix}
      x \\
      y
   \end{pmatrix}_{Q}
   +
   \begin{pmatrix}
      O_x \\
      O_y
   \end{pmatrix}_{Q} \\
   %---------------------
   \rightarrow
   \begin{pmatrix}
      x \\
      y
   \end{pmatrix}_{NE}
   &=
   \begin{pmatrix}
      \RMCOS & -\RMSIN \\
      \RMSIN & \;\;\RMCOS
   \end{pmatrix}
   \begin{pmatrix}
      x \\
      y
   \end{pmatrix}_{Q}
   +
   \begin{pmatrix}
      O_x \\
      O_y
   \end{pmatrix}_{Q} \\
   %---------------------
   \rightarrow
   \begin{pmatrix}
      x \\
      y \\
      1
   \end{pmatrix}_{NE}
   &=
   \begin{pmatrix}
      \RMCOS &    -\RMSIN & O_x \\
      \RMSIN & \;\;\RMCOS & O_y \\
           0 &          0 &   1
   \end{pmatrix}
   \begin{pmatrix}
      x \\
      y \\
      1
   \end{pmatrix}_{Q} 
   \;\;\;\; \textrm{(in augmented matrix form)}
\end{align*}


\subsection{Scale, Rotation, and Translation (SRT)}

The next case introduces two scale factors to the transformation, one for each
pixel axis in the neighboring (NW, SE, SW) sensor. These scale factors can
absorb detector tip/tilt relative to the NE sensor plane. After scale
adjustment, we expect sensors Q and NE to be coplanar.
Rotation is applied next, followed by translation. As before, $\theta$ is the
counter-clockwise rotation angle of detector Q axes w.r.t. NE detector axes. 
This model uses \textbf{5 parameters} (scales $s_1$, $s_2$, 
rotation $\theta$, and translations $O_x$, $O_y$) for each of the (NW, SE, SW)
sensors for \textbf{a total of 15 parameters}. 
\begin{align*}
   %---------------------
   \begin{pmatrix}
      x \\
      y
   \end{pmatrix}_{NE}
   &=
   \begin{pmatrix}
      \RMCOS & -\RMSIN \\
      \RMSIN & \;\;\RMCOS
   \end{pmatrix}_{Q}
   \begin{pmatrix}
      s_1 &   0  \\
        0 & s_2
   \end{pmatrix}_{Q}
   \begin{pmatrix}
      x \\
      y
   \end{pmatrix}_{Q}
   +
   \begin{pmatrix}
      O_x \\
      O_y
   \end{pmatrix}_{Q} \\
   %---------------------
   \rightarrow
   \begin{pmatrix}
      x \\
      y
   \end{pmatrix}_{NE}
   &=
   \begin{pmatrix}
      s_1 \, \RMCOS   &       -s_2 \, \RMSIN   \\
      s_1 \, \RMSIN   & \;\;\; s_2 \, \RMCOS 
   \end{pmatrix}_{Q}
   \begin{pmatrix}
      x \\
      y
   \end{pmatrix}_{Q}
   +
   \begin{pmatrix}
      O_x \\
      O_y
   \end{pmatrix}_{Q} \\
   %---------------------
   \rightarrow
   \begin{pmatrix}
      x \\
      y \\
      1
   \end{pmatrix}_{NE}
   &=
   \begin{pmatrix}
      s_1 \, \RMCOS   &     -s_2 \, \RMSIN  & O_x \\
      s_1 \, \RMSIN   & \;\; s_2 \, \RMCOS  & O_y \\
                  0   &                  0  &   1
   \end{pmatrix}_{Q}
   \begin{pmatrix}
      x \\
      y \\
      1
   \end{pmatrix}_{Q} 
   \;\;\;\; \textrm{(in augmented matrix form)}
%\]
\end{align*}

\pagebreak

\subsection{General Affine Transformation}

The final case is a general affine transformation that can account for skew,
stretch, reflection, and other things in addition to scale, rotation, and
translation. This amounts to finding a full, 4-element matrix of linear terms
(equivalent to the "CD" matrix from FITS WCS) for each sensor in addition to
the origin offset. Compared to SRT, the general affine only adds a single new
parameter but extends the model to include any affine warping (including X,Y
shear/skew and X,Y reflections that we do not expect to see). The only downside
is that the fitted parameters are products of all contributing simple
transformations and additional work is needed to factor those parameters into
contributions from the basic transformations (i.e., scale, rotation, etc).  As
before, $\theta$ is the counter-clockwise rotation angle of detector Q axes
w.r.t. NE detector axes.  This model requires \textbf{6 parameters} (4 "CD"
linear terms plus X and Y translations) for each of the (NW, SE, SW) sensors
for \textbf{a total of 18 parameters} to merge the sensors. Note that these
elements are referred to as "MCD" below to distinguish them from the true CD
matrix elements from a FITS WCS.

\begin{align*}
   %---------------------
   \rightarrow
   \begin{pmatrix}
      x \\
      y
   \end{pmatrix}_{NE}
   &=
   \begin{pmatrix}
      \textrm{M}\cdm{11} & \textrm{M}\cdm{12} \\
      \textrm{M}\cdm{21} & \textrm{M}\cdm{22} \\
   \end{pmatrix}_{Q}
   \begin{pmatrix}
      x \\
      y
   \end{pmatrix}_{Q}
   +
   \begin{pmatrix}
      O_x \\
      O_y
   \end{pmatrix}_{Q} \\
   %---------------------
   \rightarrow
   \begin{pmatrix}
      x \\
      y \\
      1
   \end{pmatrix}_{NE}
   &=
   \begin{pmatrix}
      \textrm{M}\cdm{11} & \textrm{M}\cdm{12} & O_x \\
      \textrm{M}\cdm{21} & \textrm{M}\cdm{22} & O_y \\
         0     &    0     & 1
   \end{pmatrix}_{Q}
   \begin{pmatrix}
      x \\
      y \\
      1
   \end{pmatrix}_{Q} 
   \;\;\;\; \textrm{(in augmented matrix form)}
%\]
\end{align*}
%\end{align*}

\end{document}

